Une enseigne de \qty{\v{m}}{kg} est suspendue au-dessus de l'entrée d'un magasin par deux câbles identiques. Chaque câble fait un angle de \qty{\v{th}}{\degree} avec l'horizontale, de façon symétrique. L'enseigne est immobile. Utilisez $g = \qty{9.81}{N/kg}$.

\begin{center}
\begin{tikzpicture}[thick]
    \pgfmathsetmacro{\figX}{2.6*cos(\v{th})}
    \pgfmathsetmacro{\figY}{-2.6*sin(\v{th})}
    \fill[pattern=north east lines] (-3,0) rectangle (3,0.25);
    \draw (-3,0) -- (3,0);
    \draw (-\figX,0) -- (0,\figY) -- (\figX,0);
    \draw[thin] (\figX,0) ++(180:0.7) arc (180:{180+\v{th}}:0.7);
    \path (\figX,0) ++({180+\v{th}/2}:1.1) node {\small\ang{\v{th}}};
    \draw[thin] (-\figX,0) ++(0:0.7) arc (0:{-\v{th}}:0.7);
    \draw (0,\figY) -- (0,{\figY-0.3});
    \draw[fill=gray!25] (-1.2,{\figY-0.3}) rectangle (1.2,{\figY-1.1});
    \node at (0,{\figY-0.7}) {\small\qty{\v{m}}{kg}};
\end{tikzpicture}
\end{center}

\begin{enumerate}[label=\alph*)]
    \item Déterminez le poids de l'enseigne.
    \item Déterminez la tension dans chacun des câbles.
    \item Déterminez la composante horizontale de la tension dans chaque câble. Expliquez pourquoi il serait impossible de tendre les câbles parfaitement à l'horizontale.
\end{enumerate}
